\chapter{Vitamin A (Retinol)}

\section*{What Is This Test?}
The vitamin A test measures the serum concentration of retinol, the major circulating form of vitamin A. Vitamin A is a fat-soluble vitamin that exists in several forms: preformed vitamin A (retinol, retinal, retinoic acid) found in animal sources (liver, egg yolk, dairy, fish oils), and provitamin A carotenoids (beta-carotene and others) found in orange and green vegetables, which the intestine converts to retinol. Retinol is essential for vision (retinal forms the chromophore rhodopsin in rod photoreceptors), immune function, epithelial differentiation, reproduction, and growth. The body stores approximately 80--90\% of vitamin A in the liver as retinyl esters, with only 1--2\% of total body vitamin A in the circulation, so serum retinol is a marker of circulating rather than total stores; serum retinol remains normal until hepatic stores are severely depleted. Vitamin A deficiency is a major public health problem in low-income countries, where it is the leading cause of preventable childhood blindness and is associated with increased mortality from infection. In developed countries, deficiency occurs mainly in malabsorption, liver disease, and severe dietary restriction, while hypervitaminosis A (toxicity) is seen with excessive supplementation, isotretinoin therapy, and rare liver disease.

\section*{Alternative Names}
Serum Retinol; Retinol Level; Vitamin A Level; Serum Vitamin A; Retinol Binding Protein (RBP) (related); Beta-Carotene

\section*{Principle}
Serum retinol is measured by high-performance liquid chromatography (HPLC) with ultraviolet detection at 325 nm. Serum proteins are precipitated with ethanol containing a retinol internal standard, and retinol is extracted into an organic solvent (hexane or heptane). The organic phase is dried, reconstituted, and injected onto a reverse-phase C18 column, where retinol is separated from other fat-soluble vitamins and quantified against a calibration curve. Because retinol is light-sensitive and degrades on exposure to ultraviolet light, all extraction steps are performed under subdued light or amber glassware. LC-MS/MS is an alternative method with higher throughput. Retinol binding protein (RBP), the plasma transport protein that carries retinol to tissues, is measured by immunoassay or nephelometry and correlates closely with serum retinol when nutrition is adequate; the retinol-to-RBP ratio is used in research settings to assess retinol distribution.

\section*{History}
Vitamin A was the first fat-soluble vitamin discovered. Frederick Hopkins and Elmer McCollum independently demonstrated in 1912--1913 that a fat-soluble growth factor in butter and egg yolk was essential for growth and survival; McCollum named it ``fat-soluble A,'' later designated vitamin A. In 1931, Paul Karrer determined the structure of retinol and synthesized it. George Wald's work on rhodopsin, beginning in the 1930s, explained the role of retinal in the visual cycle, earning him the 1967 Nobel Prize in Physiology or Medicine. The recognition that vitamin A deficiency causes xerophthalmia and blindness, and that supplementation dramatically reduces childhood mortality, established vitamin A supplementation programs in low-income countries. Serum retinol measurement by HPLC became clinically available in the 1970s--1980s and remains the standard laboratory assessment of vitamin A status.

\section*{Why This Test Is Ordered}
\begin{itemize}
  \item Evaluation of suspected vitamin A deficiency in patients with malabsorption syndromes (celiac disease, Crohn disease, short bowel syndrome, chronic pancreatitis, cholestatic liver disease, cystic fibrosis) who have difficulty absorbing fat-soluble vitamins
  \item Assessment of patients with alcohol-related liver disease, cirrhosis, or pancreatic insufficiency, who commonly have low vitamin A stores
  \item Evaluation of night blindness, xerophthalmia (dry eyes), Bitot spots, or other ocular manifestations of deficiency
  \item Investigation of unexplained skin changes (follicular hyperkeratosis), impaired immune function, or recurrent infections
  \item Assessment of vitamin A status in severe malnutrition or in patients on prolonged parenteral nutrition without adequate vitamin supplementation
  \item Monitoring for vitamin A toxicity in patients taking high-dose vitamin A, high-dose beta-carotene, or isotretinoin (a vitamin A derivative used for acne), especially if hepatotoxicity or intracranial hypertension is suspected
  \item Preconception counseling in at-risk populations, because deficiency and excess both have teratogenic implications
  \item Evaluation of patients with acne or other dermatologic conditions in whom long-term retinoid therapy is planned
\end{itemize}

\section*{Primary vs Secondary Test}
Serum retinol is a secondary confirmatory test, ordered when vitamin A deficiency or toxicity is clinically suspected. It is not a routine screening test in developed countries, though vitamin A screening may be part of nutritional assessment panels in malabsorption and transplant populations.

\section*{How This Test Is Performed}
A blood sample is collected by standard venipuncture into a serum separator tube. The specimen must be protected from light (wrapped in foil) because retinol is photolabile, and it should be processed promptly. Serum is analyzed by HPLC or LC-MS/MS. Retinol binding protein, when ordered, is measured by immunoassay on the same or a companion specimen. Because serum retinol is a fat-soluble vitamin measurement, it may be drawn fasting as part of a combined fat-soluble vitamin panel (A, D, E, K). Results are typically available within 3--7 days, often at a reference laboratory.

\section*{What Preparation Is Needed}
Fasting for 8--12 hours is preferred because dietary intake transiently raises serum retinol and retinyl esters. The specimen must be protected from light after collection. Patients should disclose use of vitamin A supplements, retinoids, and medications such as cholestyramine, orlistat, and mineral oil laxatives, which interfere with fat-soluble vitamin absorption and lower measured levels.

\section*{Contraindications and Precautions}
\begin{itemize}
  \item There are no absolute contraindications to standard venipuncture. Important interpretive cautions:
  \item (1) Serum retinol is not a sensitive marker of early deficiency because circulating levels are maintained until hepatic stores are nearly exhausted; low-normal values warrant interpretation in clinical context.
  \item (2) The acute-phase response lowers serum retinol: inflammation, infection, surgery, and trauma redistribute retinol and lower RBP, producing falsely low results --- measure CRP or another inflammatory marker alongside.
  \item (3) Pregnancy lowers serum retinol physiologically (hemodilution and increased demand), while breastfeeding also affects levels.
  \item (4) Oral contraceptives and estrogen therapy elevate serum retinol and RBP.
  \item (5) Renal failure elevates serum retinol and RBP.
  \item (6) Hyperlipidemia can interfere with HPLC extraction; hemolysis and exposure to light invalidate the specimen.
  \item (7) Cholestyramine, neomycin, and orlistat reduce absorption and lower serum retinol.
  \item (8) In iron-deficient or zinc-deficient patients, mobilization of hepatic retinol is impaired, so serum retinol may be low despite adequate stores.
\end{itemize}
\section*{Risks}
Minimal risk. Slight pain or bruising at the venipuncture site. Rarely: hematoma, infection, or vasovagal syncope.

\section*{What Happens After the Test}
Results are available within 3--7 days. A low serum retinol is treated with oral vitamin A supplementation (typically 10,000--25,000 IU daily for several weeks, then a lower maintenance dose), with dietary counseling emphasizing vitamin A-rich foods; ocular disease requires urgent treatment because blindness can develop rapidly. Patients with malabsorption may require water-miscible vitamin A preparations or intramuscular therapy. Vitamin A toxicity is managed by discontinuing vitamin A and retinoid supplements and avoiding vitamin A-fortified foods; symptoms (headache, blurred vision, dry skin, hair loss, bone pain, and in severe cases pseudotumor cerebri) usually resolve with removal of the source. Serum retinol is rechecked after 4--12 weeks to confirm repletion or normalization.

\section*{Understanding Results}

\subsection*{Below Normal}
\begin{itemize}
  \item \textbf{Deficiency (serum retinol below 20 mcg/dL in adults, laboratory-dependent):} Causes include inadequate dietary intake (prolonged starvation, restrictive diets), fat malabsorption (celiac disease, Crohn disease, short bowel syndrome, chronic pancreatitis, cholestasis, cystic fibrosis), alcohol-related liver disease (impaired storage and mobilization), and increased demand (pregnancy, lactation, childhood).
  \item \textbf{Clinical manifestations of deficiency:} Night blindness (earliest symptom), conjunctival xerosis, Bitot spots, corneal xerophthalmia and keratomalacia (which can progress to irreversible blindness), follicular hyperkeratosis, impaired immune function with increased severity of measles and diarrhea in children, impaired growth, and anemia.
  \item \textbf{Interpretive caution:} Serum retinol is falsely lowered by inflammation, infection, pregnancy, and hypoalbuminemia; correlate with CRP, albumin, and clinical status before diagnosing deficiency.
\end{itemize}

\subsection*{Above Normal}
\begin{itemize}
  \item \textbf{Hypervitaminosis A (serum retinol above approximately 100 mcg/dL, with elevated retinyl esters):} Caused by excessive intake of preformed vitamin A (high-dose supplements, cod liver oil), isotretinoin or other retinoid therapy, or rare hepatic disease with elevated RBP.
  \item \textbf{Acute toxicity:} Headache, dizziness, nausea, vomiting, drowsiness, and blurred vision following a single large dose.
  \item \textbf{Chronic toxicity:} Dry skin and lips, hair loss, bone and joint pain, hepatotoxicity, hypercalcemia, and pseudotumor cerebri (idiopathic intracranial hypertension with papilledema).
  \item \textbf{Teratogenicity:} Excessive vitamin A in pregnancy causes congenital malformations (neural tube and craniofacial defects), and isotretinoin is absolutely contraindicated in pregnancy.
  \item \textbf{High-dose beta-carotene} produces harmless carotenemia (yellow skin) and, in smokers, an increased risk of lung cancer in trials, so supplementation is not recommended in smokers.
  \item \textbf{Chronic renal failure} elevates serum retinol and RBP without true hypervitaminosis.
\end{itemize}

\section*{Normal Results}
\begin{itemize}
  \item \textbf{Serum retinol (adults):} 20--80 mcg/dL (0.7--2.8 $\mu$mol/L), varying by laboratory and population
  \item \textbf{Deficient:} less than 20 mcg/dL
  \item \textbf{Toxic:} greater than approximately 100 mcg/dL, especially with elevated retinyl ester fraction
  \item \textbf{Children:} Similar ranges with laboratory-specific cutoffs
  \item \textbf{Conversion:} 1 mcg/dL retinol = 0.0349 $\mu$mol/L
  \item \textbf{Daily requirement:} 900 mcg RAE (retinol activity equivalents) for adult men, 700 mcg RAE for adult women; upper limit 3,000 mcg/day preformed vitamin A
\end{itemize}

\section*{Follow-up Tests}
\begin{itemize}
  \item \textbf{Retinol binding protein (RBP):} Transport protein for retinol; low in deficiency, malnutrition, and liver disease; correlates with serum retinol when nutrition is adequate
  \item \textbf{Other fat-soluble vitamins (D, E, K):} Deficiencies commonly coexist in malabsorption and cholestasis; measure together for a complete picture
  \item \textbf{Beta-carotene:} Provitamin A status; useful when dietary intake of vegetable sources is in question
  \item \textbf{Liver function tests (ALT, AST, bilirubin, alkaline phosphatase):} Cholestatic liver disease impairs vitamin A absorption and storage
  \item \textbf{Inflammatory markers (CRP or ESR):} Interpret a low retinol in the setting of acute-phase inflammation with caution
  \item \textbf{Serum albumin and prealbumin:} Assess overall nutritional status, since hypoalbuminemia lowers RBP and retinol
  \item \textbf{Repeat serum retinol after 4--12 weeks of supplementation:} Confirms repletion and guides maintenance therapy
  \item \textbf{Ophthalmologic examination:} For evaluation and documentation of xerophthalmia or other ocular findings
\end{itemize}

\section*{References}
\begin{enumerate}
  \item Sommer A. Vitamin A deficiency and clinical disease: an historical overview. J Nutr. 2008;138(10):1835--1839.
  \item World Health Organization. Global prevalence of vitamin A deficiency in populations at risk 1995-2005. WHO; 2009.
  \item Ross AC. Vitamin A and retinoic acid in T cell-related immunity. Am J Clin Nutr. 2012;96(5):1166S--1172S.
  \item Tanumihardjo SA. Vitamin A: biomarkers of nutrition for development. Am J Clin Nutr. 2011;94(2):658S--665S.
  \item Penniston KL, Tanumihardjo SA. The acute and chronic toxic effects of vitamin A. Am J Clin Nutr. 2006;83(2):191--201.
\end{enumerate}

\clearpage
\section*{Regulatory Status and Authorization Date}
This chapter describes a test category, not a specific commercial product. FDA, CDSCO, EU IVDR, and MHRA approval or registration dates therefore cannot be assigned without the assay manufacturer, product name, instrument, and intended use. For laboratory-developed tests, an individual agency approval date may not exist. See the Preface for the interpretation of regulatory status.

\clearpage
